\documentclass[UTF8,12pt]{ctexart}
\usepackage[a4paper,margin=2.3cm]{geometry}
\usepackage{amsmath,amssymb,booktabs,graphicx,float,array}
\setlength{\parindent}{2em}
\title{波动电价下的微网日前购电模型}
\author{}
\date{}

\begin{document}
\maketitle

\begin{abstract}
问题4将外网电价由附件1的单条固定日内曲线替换为附件4的逐日逐时段波动电价，并重算问题2。本文的关键判断是：波动电价与负载、光伏一样\textbf{不可预知}，第$d$天0:00只能依据$j<d$的历史实际值预测当天144个时段的电价。因此模型在三层结构上与问题2保持一致，只把电价通道改为因果集成预测：预测电价$\hat c_{d,t}$只进入线性规划目标与储能套利，购电费用始终按当天实际电价结算。风险分位数取$q=0.80$，其依据是5倍紧急电价下的报童临界比$4c/(4c+c)$；由于计划购电与紧急购电在同一时刻按同一个$c$结算，$c$在临界比中约去，故该值与电价的取值、分布、是否可预知均无关，是事前给定的理论值，不由任何评价期回测挑选。在2025年2月1日至12月31日334天的逐日因果回放中，负载、光伏、净负荷MAE分别为148.70、149.88和247.43 kW，净负荷nMAE为8.25\%；电价预测MAE为0.0399元/kWh、nMAE为5.27\%。全年计划购电2203.21万kWh、费用1416.06万元，紧急购电7.72万kWh、费用49.35万元，总费用1465.41万元。四个指定日期中仅9月23日发生紧急购电6.807 kWh。全部48096个评价时段满足物理约束，最大能量平衡误差为$2.27\times10^{-13}$ kWh。

本文进一步用全知口径量化了波动电价的代价。同一套因果预测与同一安全余量下，若电价可预知，总费用为1460.51万元，即电价不可预知仅值4.90万元（0.33\%）。相对问题2（同一风险分位数、固定电价）高出73.58万元，其中75.2\%来自\textbf{全知最优本身就更贵}（553131.08元），仅24.8\%来自策略相对最优的距离拉大。机制是附件4的电价与负载正相关（0.5077），高负载日恰为高价日，而储能仅占日需求的17\%，无法把购电从高价日搬到低价日。

\textbf{关键词：}波动电价；微网调度；因果预测；分位数安全余量；线性规划；报童模型；全知下界
\end{abstract}

\section{问题重述}
问题4要求在电价随日期和时段波动的条件下重算问题2。除电价外，其余设定不变：微网需在每日0:00一次性确定当天144个十分钟时段的计划购电量，计划购电无论最终是否使用均按原价结算；执行时若计划购电、光伏和储能放电仍不能满足负载，不足部分按对应时段实际电价的5倍紧急采购。目标是使全年总购电费用尽可能低，并按附件5模板给出全年策略与2025年3月20日、6月21日、9月23日、12月21日的规定格式结果。

与问题2的实质差别有两点。其一是电价的\textbf{波动性}：附件4的$365\times144$价格矩阵逐日变化，而非每天重复同一条曲线。其二是电价的\textbf{不可预知性}：题面未给出电价预报，本文按与负载、光伏一致的口径处理，即当天电价只能在当天事后得知，第$d$天0:00的决策只能使用$j<d$的历史电价。

\section{问题分析}
该问题仍是带不对称损失的日前随机决策，但不确定性来源由两条通道扩展为三条：负载、光伏和电价。三条通道都必须在0:00之前用历史信息预测。

电价通道的特殊性在于它同时出现在目标函数和约束的右侧。计划购电量$x_{d,t}$由预测净负荷与安全余量决定，而费用按实际电价结算，因此电价预测误差不改变“买多少”，只改变“在哪个时段买更划算”以及储能充放电的套利方向。这一结构性差异意味着电价预测的精度要求与负载、光伏并不相同，本文在第8章对此做了定量检验。

风险分位数取$q=0.80$，与问题2保持一致，其依据是5倍紧急电价下的报童临界比$4c/(4c+c)$：由于计划购电与紧急购电在同一时刻按同一个$c$结算，$c$在临界比中约去，故该值与电价的取值、分布、是否可预知\textbf{均无关}，是事前给定的理论值（见6.1），无需任何回测。风险灵敏度网格显示最优点会随划分区间漂移（开发期0.75、留出期0.70），说明以回测选参本身缺乏稳定性；采用0.80而非开发期最优的0.75，全年费用代价为55603.82元，本文认为用可事前辩护的取值换取这一代价是合适的。

问题2与问题4采用同一风险分位数、同一套负载/光伏预测与同一安全余量，两问的对照因此只差电价通道。

\section{模型假设}
\begin{enumerate}
\item 每个十分钟时段内负载、光伏功率和电价保持不变。
\item 第$d$天0:00的预测和计划仅使用$j<d$的历史数据；负载、光伏、电价三条通道的实际值都只在日末更新。
\item 计划购电量非负，未使用的计划电量不能退款；缺口通过5倍实际电价的紧急购电补足。
\item 储能充、放电效率均为0.9，容量范围为1200--10800 kWh，充放电功率上限均为5000 kW。
\item 同一时段不同时充放电。线性规划加入微小通量惩罚，在正电价和有效率损耗下消除无意义循环。
\item 电价预测不进入光伏有效窗口与季节水平校正管线，那两个通道是光伏、负载专用的物理量。
\item 2025年1月1日0:00取$E_{1,0}=E_0=6000$ kWh，此后满足跨日连续$E_{d+1,0}=E_{d,T}$；一月参与预测预热与状态递推，正式费用评价期为2月1日至12月31日。
\item 第0天的电价冷启动先验取附件1的单条曲线，与负载、光伏的initial先验同构。
\item 不使用附件3的光伏预报；那是问题3的机制，本问与问题2保持一致。
\end{enumerate}

\section{符号说明}
\begin{table}[H]\centering\small
\begin{tabular}{lll}\toprule
符号 & 含义 & 单位\\\midrule
$d$ & 日期序号，$d=1,\ldots,D$ & ---\\
$t$ & 日内十分钟时段，$t=1,\ldots,T$ & ---\\
$D,T$ & 全年天数与每日时段数，$D=365,T=144$ & ---\\
$L_{d,t},P_{d,t}^{\mathrm{act}}$ & 小区负载功率、光伏实际功率 & kW\\
$L_{d,t}^{\mathrm{fc}},P_{d,t}^{\mathrm{fc}}$ & 日前负载预测、光伏功率预测 & kW\\
$c_{d,t}$ & 当天逐时段\textbf{实际}外网电价（结算用） & 元/kWh\\
$\hat c_{d,t}$ & 第$d$天0:00可用的电价预测（仅入LP目标） & 元/kWh\\
$\bar c$ & 电价均值 & 元/kWh\\
$m_{d,t}$ & 净负荷分位数安全余量 & kW\\
$q$ & 安全余量分位数 & ---\\
$x_{d,t}$ & 外网计划购电功率（全额计价） & kW\\
$y_{d,t}$ & 实际使用的计划购电功率，$0\le y_{d,t}\le x_{d,t}$ & kW\\
$r_{d,t}=x_{d,t}-y_{d,t}$ & 购电浪费功率 & kW\\
$g_{d,t}$ & 光伏实际使用功率 & kW\\
$w_{d,t}=P_{d,t}^{\mathrm{act}}-g_{d,t}$ & 光伏浪费功率 & kW\\
$u_{d,t},v_{d,t}$ & 储能充电功率、放电功率 & kW\\
$E_{d,t}$ & 时段末储电量 & kWh\\
$E_0,E_{\min},E_{\max}$ & 初始、最小和最大储电量 & kWh\\
$P_{\mathrm{st}}^{\max}$ & 最大充/放电功率 & kW\\
$e_{d,t}$ & 外网紧急购电功率 & kW\\
$\eta_c,\eta_d$ & 充、放电效率 & ---\\
$\Delta t$ & 时段长度，$1/6$小时 & h\\\bottomrule
\end{tabular}
\end{table}

\section{因果预测模型}
\subsection{负载与光伏}
负载与光伏预测沿用问题2最终版的季节自适应集成，未作任何改动：基础层由近期加权基线、岭回归和浅层LightGBM构成并在线加权，校正层用历史光伏有效窗口和日能量比作季节水平修正。基础模型的近期损失为
\[
s_{h,d}=0.85\,\overline \ell_{h,d}^{\mathrm{EW}}+0.15Q_{0.75}(\ell_{h,d}),
\]
校正后的预测为
\[
L_{d,t}^{\mathrm{fc}}=\lambda^{L}_dL_{d,t}^{\mathrm{fc},(0)},\qquad
P_{d,t}^{\mathrm{fc}}=\lambda^{P}_dP_{d,t}^{\mathrm{fc},(0)}\mathbf 1(t\in W_d).
\]
消融实验已表明周效应通道会使净负荷MAE与总费用同时变差，故关闭。

\subsection{电价的因果集成预测}
电价使用与负载、光伏同构的三模型因果集成：近期加权基线、岭回归和浅层LightGBM，按近期损失自适应加权，训练窗随残差漂移缩放。第$d$天的特征只包含$j<d$的滞后值、同时段近期均值、日内周期编码、星期编码与年内日序编码；第0天用附件1的单条曲线作冷启动先验。LightGBM采用固定随机种子、单线程与确定性模式，保证重跑逐位一致。
\[
\hat c_{d,t}=\sum_{h}\omega_{h,d}\,\hat c^{(h)}_{d,t},\qquad
\omega_{h,d}\propto \exp(-\eta\, s_{h,d}).
\]
需要强调的是，电价预测\textbf{不进入}光伏有效窗口与季节水平校正管线：那两个通道依赖的是光伏出力的物理时段与日电量水平，对电价没有意义。

\subsection{预测结果}
\begin{table}[H]\centering
\caption{正式评价期预测精度}
\begin{tabular}{lrrrr}\toprule
目标 & MAE & RMSE & nMAE & 偏差\\\midrule
负载/kW & 148.70 & 197.72 & 3.22\% & 2.10\\
光伏/kW & 149.88 & 304.11 & 6.30\% & 3.09\\
净负荷/kW & 247.43 & 362.77 & 8.25\% & $-0.99$\\
电价/(元/kWh) & 0.0399 & 0.0540 & 5.27\% & $+0.0013$\\\bottomrule
\end{tabular}
\end{table}

电价的实际均值为0.7575元/kWh（正式评价期），预测偏差仅$+0.0013$元/kWh，说明集成的水平无系统性偏移。

\begin{figure}[H]\centering
\includegraphics[width=0.95\textwidth]{../figures/电价预测.png}
\caption{附件4电价的因果预测与实际值对比}
\end{figure}

\section{风险控制与调度模型}
\subsection{报童临界比：$q=0.80$与电价无关}
设某时段的实际净负荷为$N$（随机），计划购电量为$x$。该时段费用为
\[
C(x)=c\,x+5c\,(N-x)^+,
\]
其中$c$为\textbf{该时段的实际电价}。对$x$求一阶条件：
\[
\frac{\partial}{\partial x}E[C]=E[c]-5E\!\left[c\,\mathbf 1(N>x)\right]
=E[c]-5P(N>x)E[c\mid N>x]=0,
\]
即
\[
P(N>x)\,E[c\mid N>x]=\frac{E[c]}{5}.
\]
关键在于：计划购电与紧急购电在\textbf{同一时刻按同一个$c$结算}，$c$在等式中被约去，临界比退化为常数
\[
\frac{4c}{4c+c}=0.80 .
\]
因此$q=0.80$与电价的取值、分布、是否可预知\textbf{均无关}，也不依赖负载与电价的独立性假设，是事前给定的理论值。这与问题2的情形相比是一个更强的论证：问题2的0.75来自开发期回测，而本问的0.80由模型结构直接给出，不存在样本内选择问题。

安全余量的构造沿用问题2的长短窗混合分位数并作平滑与阻尼，第$d$天只使用$j<d$的净负荷残差：
\[
\widetilde m_{d,t}=0.7Q_q(\xi_{d-90:d-1,t})+0.3Q_q(\xi_{d-30:d-1,t}),\qquad
m_{d,t}=0.35\widetilde m_{d,t}+0.65m_{d-1,t}.
\]
计划净需求为$N^{\mathrm{plan}}_{d,t}=N_{d,t}^{\mathrm{fc}}+m_{d,t}$。需要注意安全余量只作用于净负荷，与电价通道无关，因此问题2与问题4的余量完全相同，两问的对照可以干净地归因到电价。

\subsection{日前线性规划}
在每一天给定当日初始储电量后，求解
\[
\min \sum_{t=1}^{T}\hat c_{d,t}\,x_{d,t}\Delta t-\lambda_E E_{d,T}
+\varepsilon_0\sum_t(u_{d,t}+v_{d,t}+w_{d,t})\Delta t,
\]
其中$\lambda_E=\eta_d\,\bar{\hat c}_d$为日末储能水值，按当天\textbf{预测}电价的均值计，$\varepsilon_0$为消除无意义循环流的微小正数。计划阶段约束为
\[
x_{d,t}+P_{d,t}^{\mathrm{fc}}-w_{d,t}+v_{d,t}
=L_{d,t}^{\mathrm{fc}}+m_{d,t}+u_{d,t},
\]
\[
E_{d,t}=E_{d,t-1}+\eta_cu_{d,t}\Delta t-\frac{v_{d,t}\Delta t}{\eta_d},
\]
\[
E_{\min}\le E_{d,t}\le E_{\max},\qquad
0\le u_{d,t},v_{d,t}\le P_{\mathrm{st}}^{\max},\qquad
x_{d,t},g_{d,t},w_{d,t}\ge0 .
\]
执行阶段使用实际值，并统一符号关系
\[
y_{d,t}+g_{d,t}+v_{d,t}+e_{d,t}=L_{d,t}+u_{d,t},\qquad
g_{d,t}+w_{d,t}=P_{d,t}^{\mathrm{act}},
\]
其中$r_{d,t}=x_{d,t}-y_{d,t}\ge0$为购电浪费功率，$e_{d,t}$为紧急购电功率。\textbf{结算始终使用实际电价$c_{d,t}$}，全年实际总费用为
\[
C=\sum_{d=32}^{D}\sum_{t=1}^{T}
\left(c_{d,t}x_{d,t}+5c_{d,t}e_{d,t}\right)\Delta t .
\]
对应地，$C_d^{\mathrm{grid}}=\sum_t c_{d,t}x_{d,t}\Delta t$，$C_d^{\mathrm{emer}}=\sum_t5c_{d,t}e_{d,t}\Delta t$，且$C_d=C_d^{\mathrm{grid}}+C_d^{\mathrm{emer}}$。注意目标函数用预测电价$\hat c$、结算用实际电价$c$，这是电价不可预知性的唯一体现；若令$\hat c=c$即退化为完美电价信息的对照口径。程序将功率乘以$\Delta t$后保存为每十分钟电量，因此结果CSV中的购电、充放电和紧急购电单位均为kWh。

\section{结果}
\subsection{全年结果}
\begin{table}[H]\centering
\caption{2025年2月1日至12月31日（334天）结果}
\begin{tabular}{lr}\toprule
指标 & 数值\\\midrule
计划购电量/kWh & 22032122.53\\
计划购电费/元 & 14160574.10\\
紧急购电量/kWh & 77165.86\\
紧急购电费/元 & 493503.49\\
总费用/元 & \textbf{14654077.59}\\
未用计划电量/kWh & 1435523.93\\
弃光量/kWh & 1432173.62\\
发生紧急购电的日期数 & 112\\
单日费用95\%分位数/元 & 74434.16\\
单日费用CVaR95/元 & 80780.77\\
期末储电量/kWh & 5924.85\\\bottomrule
\end{tabular}
\end{table}

计划购电加权均价为0.6427元/kWh，明显低于同期实际均价0.7575元/kWh，说明模型在低价时段与低价日主动多购。紧急购电的加权单价为6.395元/kWh，约为计划购电均价的10倍；这是紧急购电集中于电价高峰时段的必然结果（问题2在固定电价下同样呈现10.23倍的比值），不是问题4引入的偏差。

\subsection{指定日期表1结果}
单位：各时段和全天购电量为kWh，费用为元。
\begin{table}[H]\centering\small
\caption{表1格式的计划购电结果}
\resizebox{\textwidth}{!}{\begin{tabular}{lrrrrrrrr}\toprule
日期&10:00&12:00&14:00&16:00&18:00&20:00&全天购电量&全天购电费\\\midrule
3月20日&0.00&625.33&0.00&494.00&383.42&759.32&69764.55&43223.00\\
6月21日&0.00&0.00&0.00&238.56&478.01&0.00&38130.30&19086.46\\
9月23日&0.00&534.14&0.00&605.30&0.00&0.00&73786.83&48263.63\\
12月21日&0.00&974.72&0.00&830.86&724.94&0.00&97399.86&73243.94\\\bottomrule
\end{tabular}}
\end{table}

\subsection{指定日期表2结果}
表中“充/放”为对应4小时内累计充电量/放电量，单位kWh。
\begin{table}[H]\centering\small
\caption{表2格式的储能结果}
\resizebox{\textwidth}{!}{\begin{tabular}{lrrrrrrrr}\toprule
日期&0--4&4--8&8--12&12--16&16--20&20--24&0:00储电&24:00储电\\\midrule
3月20日&88.51/108.15&335.59/5571.95&6285.65/2777.38&4046.04/450.44&258.70/6993.70&10366.45/1416.95&10800.00&10800.00\\
6月21日&264.14/7329.99&1246.57/1737.46&9683.67/0.00&0.00/0.00&0.00/5451.30&9799.01/2485.90&10800.00&10800.00\\
9月23日&0.28/0.23&312.19/5274.70&6293.08/1896.62&2781.88/432.27&154.69/5838.56&10121.33/2485.02&10800.00&10800.00\\
12月21日&55.38/134.20&136.80/842.52&5602.28/7281.15&7145.60/2226.67&24.51/5743.00&10578.73/2842.53&10800.00&10800.00\\\bottomrule
\end{tabular}}
\end{table}

\subsection{指定日期紧急购电}
\begin{table}[H]\centering
\caption{指定日期紧急购电量}
\begin{tabular}{lrrr}\toprule
日期&紧急购电/kWh&紧急购电费/元&全天总费用/元\\\midrule
2025-03-20&0.000&0.00&43223.00\\
2025-06-21&0.000&0.00&19086.46\\
2025-09-23&6.807&44.08&48307.71\\
2025-12-21&0.000&0.00&73243.94\\\bottomrule
\end{tabular}
\end{table}

\begin{figure}[H]\centering
\includegraphics[width=0.95\textwidth]{../figures/指定日期调度.png}
\caption{四个指定日期的负载、光伏及购电策略}
\end{figure}

\section{波动电价与不可预知性的代价}
本章的全部对照口径都使用\textbf{同一套因果负载/光伏预测与同一$q=0.80$安全余量}，只替换电价通道，因此差额可以干净地归因到电价。所有“全知”口径都不属于可实现的策略，仅用于量化代价，本文的可实现策略始终严格遵守$j<d$的信息集。

\subsection{代价阶梯}
\begin{table}[H]\centering
\caption{四个口径的总费用}
\begin{tabular}{lrr}\toprule
口径 & 总费用/元 & 相对全知下界\\\midrule
理论完美最低值（全知、全年统一优化） & 12782591.81 & ---\\
完美预测但逐日决策（余量0） & 12939133.66 & $+156541.85$（$+1.22\%$）\\
完美电价、负载/光伏仍因果预测 & 14605084.97 & $+1822493.16$（$+14.26\%$）\\
可实现策略（本文提交） & 14654077.59 & $+1871485.78$（$+14.64\%$）\\\bottomrule
\end{tabular}
\end{table}

“理论完美最低值”是一次性求解全年$365\times144=52560$个时段的统一线性规划，已知全部实际负载、光伏与电价，储能约束、初始SOC$=6000$ kWh、年末SOC回归6000 kWh与主模型完全一致，费用仅统计正式评价期。它是任何因果策略都不可能低于的下界。“完美预测但逐日决策”则把三条通道的预测都换成实际值、安全余量置0，但仍保持每天0:00只解当天144个时段、日末储电按水值$\lambda_E=\eta_d\bar c$折价的逐日结构。

两个全知口径相差156541.85元（1.22\%），衡量的是“逐日滚动$+$水值近似”这一决策结构的次优性；由于两者都不含任何预测误差，这1.22\%与预测精度无关。\textbf{换言之，主模型的调度结构本身几乎没有浪费，14.64\%的差距中有13.4\%来自预测。}

\subsection{与问题2的费用差分解}
为与问题2对照，把电价换回附件1的单条固定曲线重算，得到13918258.75元，与问题2的正式程序结果逐位一致，说明两问的管线只差电价通道。

\begin{table}[H]\centering\small
\caption{问题2与问题4的三层口径对照}
\resizebox{\textwidth}{!}{\begin{tabular}{lrrrr}\toprule
口径 & 全知全年统一下界/元 & 完美预测逐日/元 & 可实现/元 & 可实现$-$下界\\\midrule
问题2（附件1固定电价） & 12229460.73 & 12288748.57 & 13918258.75 & 1688798.03（13.81\%）\\
问题4（附件4波动电价） & 12782591.81 & 12939133.66 & 14654077.59 & 1871485.78（14.64\%）\\\bottomrule
\end{tabular}}
\end{table}

两问费用差735818.83元的分解如下（逐项相加与总差完全相等）：
\begin{table}[H]\centering
\caption{费用差的来源分解}
\begin{tabular}{lrr}\toprule
来源 & 金额/元 & 占比\\\midrule
① 全知最优本身抬高 & 553131.08 & 75.2\%\\
② 逐日决策结构的次优性变大 & 97254.01 & 13.2\%\\
③ 预测误差与安全余量的代价变大 & 85433.75 & 11.6\%\\\bottomrule
\end{tabular}
\end{table}

需要指出，两问相对下界的\textbf{比例}只差0.83个百分点（13.81\%$\to$14.64\%），差别主要不在“策略做得差”，而在“地板被抬高”：
\begin{itemize}
\item \textbf{①（75.2\%）}是物理性劣势。问题2的电价是单一固定曲线，每天形状相同，全知下模型可以把购电精确排进该曲线的低价时段；问题4的电价与负载正相关（0.5077），高负载日恰为高价日，而储能仅占日需求的17\%，即使全知也无法把电从高价日搬到低价日。这部分任何策略都躲不掉。
\item \textbf{②（13.2\%）}来自日末水值近似。问题2每天电价相同，$\lambda_E=\eta_d\bar c$很准；问题4每天电价不同，用当天均价近似“把电留到明天”的真实机会成本误差更大，逐日LP的充放电取舍偏差被放大，逐日结构的损失由5.93万元涨到15.65万元。
\item \textbf{③（11.6\%）}来自误差穿透余量后的紧急购电。安全余量与电价无关，但问题4中紧急购电按波动的实际电价乘5倍结算，惩罚更贵，代价由162.95万元涨到171.49万元。
\end{itemize}

\subsection{机制：电价与负载同向波动}
进一步把可实现的电价通道拆成三步，定位差额究竟来自日内形状、日间波动还是不可预知：
\begin{table}[H]\centering
\caption{电价通道的三步分解}
\begin{tabular}{lrrr}\toprule
口径 & 总费用/元 & 增量 & 占比\\\midrule
C1 附件1单条曲线（已知） & 13918258.75 & --- & ---\\
C2 附件4日内平均形状（已知，无逐日波动） & 13918261.47 & $+2.72$ & 0.0\%\\
C3 附件4实际电价（完美可预知） & 14605084.97 & $+686823.49$ & 93.3\%\\
C4 附件4实际电价（因果预测） & 14654077.59 & $+48992.62$ & 6.7\%\\\bottomrule
\end{tabular}
\end{table}

结论有三点：
\begin{enumerate}
\item \textbf{日内形状没有差异。}附件4按全期时段平均后与附件1曲线几乎完全重合（相关系数0.9997，全期均值均为0.7662，最低价0.3713、最高价1.3952完全相同；按正式期平均则为均值0.7575、最低0.3660、最高1.3824），因此C2与C1只差2.72元。问题4与问题2的全部差异都在\textbf{日与日之间}。
\item \textbf{波动本身占93.3\%，不可预知只占6.7\%。}即使电价完全可预知，也要多付686823.49元；而“电价不可预知”只值48992.62元。
\item \textbf{购电量几乎不变。}22011526.88$\to$22032122.53 kWh（$+0.09\%$），因为购电量由净负荷与安全余量决定，储能只能微调。
\end{enumerate}

波动之所以推高成本而不是提供套利机会，根源是\textbf{电价与负载正相关（0.5077）}：日购电量与日均价的相关系数高达0.9646。电价最高的1/3日子平均日购电85779 kWh、日均价0.8582元/kWh，电价最低的1/3日子为41737 kWh、0.6376元/kWh——买得多的日子电价贵34.6\%，购电量还是后者的2.06倍。日内择时能力同时被削弱：购电量与电价的相关系数由问题2的$-0.5015$退化到问题4的$-0.3413$。

也就是说，波动给出的“低价日囤电”机会因储能太小（10.8 MWh vs 日需求约65 MWh）而用不上，风险却躲不掉。一个直接证据是：把问题2的购电量分布\textbf{原封不动}乘上附件4实际电价，计划购电费即由13418699.15元升至14158188.07元（加权均价0.6096$\to$0.6432，$+739488.92$元），这一项就解释了几乎全部差额；问题4模型重新优化后还少花约3670元，说明策略本身并未变差。

\subsection{安全余量的成本与收益}
由于电价在两个对照口径中都是完美的，8.1中“完美电价”与“完美预测”之间的1665951.31元全部来自负载/光伏侧。将预测与余量交叉组合可分解如下：
\begin{table}[H]\centering\small
\caption{预测误差与安全余量的四口径交叉}
\resizebox{\textwidth}{!}{\begin{tabular}{lrrrrr}\toprule
口径（电价均为附件4实际值） & 总费用/元 & 计划购电量/kWh & 紧急购电量/kWh & 紧急天数 & 未用计划电量/kWh\\\midrule
W1 实际值 $+$ 余量0 & 12939133.66 & 20386278 & 0 & 0 & 0\\
W2 实际值 $+$ $q=0.80$ & 14138404.74 & 22123188 & 0 & 0 & 1356143\\
W3 因果预测 $+$ 余量0 & 15891421.11 & 20389784 & 465765 & 304 & 398871\\
W4 因果预测 $+$ $q=0.80$ & 14605084.97 & 22096678 & 79894 & 123 & 1442835\\\bottomrule
\end{tabular}}
\end{table}

线性分解会出现$-2485607.23$元的交互项，说明预测误差与安全余量强耦合、不可加（余量本质是对误差的保险），因此改用两条顺序分解，二者都等于1665951.31元：
\begin{itemize}
\item 先余量、后误差：W1$\to$W2为$+1199271.08$，W2$\to$W4为$+466680.23$；
\item 先误差、后余量：W1$\to$W3为$+2952287.46$，W3$\to$W4为$-1286336.15$。
\end{itemize}

由此得到两点结论。其一，\textbf{安全余量单独看是纯浪费}：若预测完美仍留$q=0.80$的余量，会多买1356143 kWh用不完的电（占购电6.1\%），白花119.93万元。其二，\textbf{余量在真实预测下是划算的保险}：不留余量直接按预测采购，误差全部转化为5倍价紧急购电（465765 kWh、304天、紧急费298.62万元），总费用升至15891421元；留余量后紧急购电压至79894 kWh / 123天 / 51.90万元，\textbf{净省128.63万元}。因此在当前预测精度下，$q=0.80$是正收益的选择。

与此相关的一个细节是购电浪费落在什么价位：未用的1435524 kWh加权均价仅0.5765元/kWh，低于计划购电均价0.6427元/kWh；$\mathrm{corr}$(未用电量,电价)$=-0.1572$，电价最高20\%时段的浪费率只有1.73\%，最低20\%时段为6.79\%。即模型在低价时段主动多买，多买的部分自然最容易用不完——浪费是“便宜时买多了”，而不是“买贵了”。

\section{灵敏度与稳健性分析}
\subsection{风险分位数灵敏度}
\begin{table}[H]\centering
\caption{风险分位数灵敏度（仅事后对照，不参与选参）}
\begin{tabular}{rrrrrr}\toprule
$q$&开发期费用/元&留出期费用/元&全年费用/元&全年紧急购电/kWh&紧急购电天数\\\midrule
0.60&9387547.47&5612149.26&14999696.73&263865.78&246\\
0.65&9188012.97&5555222.07&14743235.05&194017.74&211\\
0.70&9078866.51&\textbf{5548003.43}&14626869.94&145058.12&177\\
0.75&\textbf{9028790.49}&5569683.27&\textbf{14598473.77}&106860.02&139\\
\textbf{0.80}（采用）&9045621.28&5608456.31&14654077.59&77165.86&112\\
0.85&9125937.17&5680493.13&14806430.30&53548.60&80\\\bottomrule
\end{tabular}
\end{table}

该网格\textbf{仅作事后对照，不参与选参}。开发期最优为0.75、留出期最优为0.70，与问题2完全相同的漂移模式，说明“最优点随划分区间移动”是模型固有性质，与电价是否波动无关。这也正是本文改用事前给定的$q=0.80$的理由：全年费用在0.75处最低（14598473.77元），采用0.80的代价为55603.82元，但换来了不依赖任何评价期回测的可辩护性。

\subsection{物理可行性}
全年48096个正式评价时段中，储电量越界、充放电功率越界、同时充放电和负流量次数均为0；最大能量平衡误差为$2.27\times10^{-13}$ kWh。因此数值解满足全部物理约束。

\subsection{结论边界}
数据只覆盖一个自然年度，无法验证跨年度泛化，尤其无法判断“电价与负载正相关”是否为长期性质；若该相关性在未来减弱，问题4相对问题2的劣势会相应缩小。模型没有使用电价与气象的外生预报，对极端价格尖峰的预测能力有限。此外，第8章的全知口径只用于量化代价，其中“年末SOC回归6000 kWh”的约束会抬高下界，若放开该约束，理论完美最低值会更低。

\section{模型评价}
模型的优点在于：三条通道严格因果，信息集与题面一致；风险分位数由报童临界比在电价未知条件下推出，不依赖任何回测；LP目标与结算的电价分离，使“不可预知性”可以被干净地度量；全部结果与全部对照口径都由同一主程序与诊断脚本复现。

不足在于：电价预测目前是跟随型集成，未利用电价与负载的跨通道相关性（0.5077），若把负载预测作为特征或对两者残差作联合建模，日内择时能力（当前相关系数$-0.3413$，低于问题2的$-0.5015$）仍有改善空间；安全余量仍是全局分位数，未按季节或电价高低分档；单年度数据不能证明“电价—负载正相关”的长期稳定性。

\section{结论}
问题4最终采用“负载/光伏季节自适应预测$+$电价因果三模型集成$+$0.80报童分位数安全余量$+$确定性储能LP”的结构，预测电价只进入LP目标与储能套利，费用按实际电价结算。全年总费用为14654077.59元，四日结果与全年逐十分钟策略均由同一模型直接生成，\texttt{results/result4-2.xlsx}为对应提交结果。

本问最重要的结论不是费用本身，而是费用的归因：问题4比问题2贵的73.58万元中，\textbf{75.2\%来自波动电价使全知最优本身抬高55.31万元}（电价与负载同向波动、储能太小），只有13.2\%来自逐日水值近似、11.6\%来自预测误差。而“电价不可预知”这一最容易被归咎的因素，实测只值4.90万元（0.33\%）。因此针对本问的改进方向应是提高储能的跨日搬运能力或在低价日增大计划购电的容错，而非继续压低电价预测的MAE。

\end{document}
